\begin{abstract}
Unified multimodal models for dense grounded understanding must support both language interaction and pixel-level outputs over images and videos, which leads to large visual token budgets and high inference cost.
While learnable visual token selection has shown promise in VQA-style multimodal large language models, it remains unclear whether it still works in unified dense grounding systems, where token reduction may directly affect segmentation-sensitive intermediate representations.
In this work, we study this question on Sa2VA and present \textbf{Sa2VA-Select}, a compatible extension that transfers learnable visual token selection to Sa2VA's unified visual stream before the multimodal LLM.
Rather than redesigning the baseline task interface, Sa2VA-Select preserves Sa2VA's original segmentation-and-conversation pipeline and learns which visual tokens should be retained under a given budget.

Experiments show that learnable token selection remains effective in this more demanding setting.
On Sa2VA-1B, Sa2VA-Select preserves strong multi-task performance under substantial token reduction and improves several dense grounding benchmarks at moderate retention ratios.
At 30\% retention, it improves ReVOS from 47.6 to 51.3 while reducing peak GPU memory from 2.76 to 2.46 GB and prefill time from 40.16 to 35.14 ms.
Across retention ratios, 30--40\% emerges as the preferred operating region, offering the most favorable trade-off between dense grounding quality and efficiency in our current setting.
On Ref-SAV, our 1B model reaches 57.6 J\&F, surpassing both Sa2VA-8B (50.0) and zero-shot Sa2VA-8B (41.3) under our evaluation protocol.
These results show that learnable token selection can be compatibly integrated into unified dense grounded image-video models and provide a favorable accuracy--efficiency trade-off.
\end{abstract}
